\section{Related Work}
\label{sec:related}

\subsection{Structure-Based Visual Localization}
Structure-based localization estimates camera pose from image-to-3D correspondences~\cite{hartley2004mvg,lepetit2009epnp,fischler1981ransac}. Classical pipelines obtain reference geometry from structure from motion (SfM), where matched image observations are organized into tracks and triangulated into 3D points~\cite{schonberger2016sfm}. Learned local features and matchers have substantially improved the robustness of this correspondence stage~\cite{detone2018superpoint,sarlin2020superglue,sarlin2019hloc}. In parallel, pose-regression and scene-coordinate methods learn a direct mapping from image evidence to camera pose or scene coordinates~\cite{kendall2015posenet,shotton2013scenes,brachmann2021dsacstar,brachmann2023ace,reported260504730v1bibbib37}. Our method remains in the explicit correspondence-based family; its distinction is that the localization map is derived from an already reconstructed Gaussian scene rather than built from the original mapping images.

Compact map construction is a related concern because an explicit map must remain searchable as its landmark count grows. Prior work has studied observation prioritization, image coverage, and compact representations~\cite{li2010prioritized,cao2014minimal,dymczyk2015gist}, as well as nonuniform appearance allocation and pose-aware map compression~\cite{camposeco2019hybrid,yang2022scenesqueezer}. These methods motivate selecting landmarks according to their contribution to localization rather than retaining all available structure. Our selection operates on landmarks reconstructed from rendered tracks and explicitly combines reference-view support with local pose information.

\subsection{Localization from Renderable Scene Representations}
Neural scene representations can also serve as sources of localization observations. For NeRFs~\cite{mildenhall2020nerf}, iNeRF performs pose estimation by minimizing a rendering-based photometric objective~\cite{yenchen2021inerf}; subsequent feature-based methods use learned correspondences or rendered observations to make alignment more robust~\cite{reported260504730v1bibbib42,reported260504730v1bibbib27,zhou2024nerfmatch,reported250319358v1bibbib68,reported260504730v1bibbib57}. Renderable models can likewise provide offline reference images or training data for localization~\cite{reported260504730v1bibbib56,reported260504730v1bibbib29,reported260504730v1bibbib8}.

MeshLoc~\cite{panek2022meshloc} is particularly close in spirit: it uses a renderable mesh to obtain reference-view depth and 3D correspondences, while also studying triangulation from matched reference features. Our conversion instead builds a persistent landmark map offline from feature tracks between rendered views. Once constructed, the map no longer requires access to the renderable scene, and each landmark stores a single fused descriptor rather than a reference image or a rendering primitive.

\subsection{Gaussian-Based Localization}
Recent Gaussian-based localizers use the scene representation itself to support feature matching and pose estimation. SplatLoc learns a scene-specific descriptor decoder and selects Gaussian primitives~\cite{zhai2025splatloc}. STDLoc learns feature-bearing Gaussians and a detector, followed by sparse pose initialization and dense refinement~\cite{huang2025stdloc}. Other approaches transfer descriptors to Gaussian representations or use Gaussian-assisted refinement~\cite{reported260504730v1bibbib47,reported260504730v1bibbib52,reported260504730v1bibbib61}. LiteLoc addresses representation size while retaining a sparse-to-dense localization procedure~\cite{li2026liteloc}.

ULF-Loc instead replaces blended feature optimization with geometry-weighted multiview fusion associated with Gaussian landmarks~\cite{gu2026ulfloc}. SplitGS-Loc reduces descriptor ambiguity by splitting and selecting primitives, with an enhanced variant using Gaussian-assisted refinement~\cite{lee2026splitgsloc}. Other methods localize against prebuilt Gaussian scenes using generalizable encoders or rendered-feature matching~\cite{khatib2025gsvisloc,lee2025gsfeatloc}. These methods demonstrate the value of Gaussian scenes for localization, but they generally retain the rendering representation as the localization substrate. Our approach uses Gaussian rendering as an offline observation source and exports a separate explicit landmark map whose coordinates are determined by multiview feature correspondences. This separation enables the same map to be queried without the Gaussian renderer.